% ============================================================================
% Trang đầu theo format HNUE chính thức:
%   - 2 KHỐI DỌC: TIẾNG VIỆT trước, TIẾNG ANH sau (KHÔNG phải 2 cột ngang).
%   - Title 11pt UPPERCASE BOLD centered; cách dòng giữa các phần.
%   - Author 11pt; affiliation 11pt italic.
%   - *Tác giả liên hệ + e-mail 11pt; ngày tháng 10pt.
%   - Tóm tắt + từ khoá 10pt, lề ±1cm, không indent đầu đoạn.
% ============================================================================

\enlargethispage{1.5cm}

% ─── KHỐI TIẾNG VIỆT ──────────────────────────────────────────────────────
\begin{center}
{\fontsize{11}{13}\bfseries\selectfont
ƯỚC LƯỢNG NĂNG LỰC CỦA HỌC SINH THEO THỜI GIAN THỰC DỰA TRÊN LÝ THUYẾT ĐÁP ỨNG CÂU HỎI (IRT) VÀ MÔ HÌNH PHƯƠNG TRÌNH CẤU TRÚC DỌC (LSEM)}
\end{center}

\vspace{6pt}

\begin{center}
{\fontsize{11}{13}\selectfont
Nguyễn Tiến Đạt\textsuperscript{1,\,*}, Nguyễn Diệp Linh\textsuperscript{1}, Bạch Đức Anh\textsuperscript{1}, Bùi Thị Thu Thuỷ\textsuperscript{1} và Phạm Lê Nhật Linh\textsuperscript{1}}
\end{center}

\begin{center}
{\fontsize{11}{13}\itshape\selectfont
\textsuperscript{1}Khoa Công nghệ Thông tin, Trường Toán học và Công nghệ thông tin,\\
Trường Đại học Sư phạm Hà Nội, Hà Nội, Việt Nam}
\end{center}

\begin{center}
{\fontsize{11}{13}\selectfont
\textsuperscript{*}Tác giả liên hệ: Nguyễn Tiến Đạt, e-mail: \href{mailto:mas835151465@hnue.edu.vn}{mas835151465@hnue.edu.vn} (chính); \href{mailto:datlcpro@gmail.com}{datlcpro@gmail.com} (phụ)}\\[1pt]
{\fontsize{10}{12}\selectfont
Ngày nhận bài: \today. Ngày sửa bài: TBD. Ngày nhận đăng: TBD.}
\end{center}

\vspace{4pt}

{\fontsize{10}{12}\selectfont
\setlength{\leftskip}{1cm}\setlength{\rightskip}{1cm}
\noindent\textbf{Tóm tắt.} Các hệ thống Lý thuyết Đáp ứng câu hỏi (IRT) triển khai quy mô lớn, kể cả nền tảng OLM của Việt Nam, thường giả định năng lực học sinh \(\theta\) là tĩnh, trong khi dữ liệu học trực tuyến dày đặc theo thời gian mâu thuẫn với giả định đó. Bài báo xây dựng và kiểm chứng khung tích hợp IRT--LSEM ước lượng năng lực động \(\theta_{it}\) trên dữ liệu Toán THPT từ OLM gồm 24{.}213 học sinh và 122{.}957 lượt đánh giá khối 10--12. Khung gồm tầng đo lường 1PL (EAP) kết hợp bốn tầng cấu trúc dọc (LGCM, LCSM, mô hình đa cấp, CT-DSEM) và bộ lọc Kalman dùng sai số EAP làm phương sai quan sát biến thiên. Mô hình động vượt baseline tĩnh với \(\Delta\)AIC = 82--167 (LRT \(p \le 10^{-19}\)); LCSM phát hiện cơ chế đuổi kịp ở khối 10 (\(\beta = -0{,}069\), CI loại trừ 0) chuyển sang trạng thái ổn định (\(\beta \approx 0\)) ở khối 12; Kalman giảm 32--41\% phương sai sai số chuẩn EAP mà không cần thêm câu hỏi.\par
\noindent\textbf{\textit{Từ khoá:}} \textit{Lý thuyết đáp ứng câu hỏi, mô hình phương trình cấu trúc dọc, bộ lọc Kalman, ước lượng năng lực động, khai phá dữ liệu giáo dục.}
\par}

\vspace{6pt}

% ─── KHỐI TIẾNG ANH ──────────────────────────────────────────────────────
\begin{center}
{\fontsize{11}{13}\bfseries\selectfont
REAL-TIME ESTIMATION OF STUDENT ABILITY BASED ON ITEM RESPONSE THEORY (IRT) AND LONGITUDINAL STRUCTURAL EQUATION MODELING (LSEM)}
\end{center}

\vspace{6pt}

\begin{center}
{\fontsize{11}{13}\selectfont
Nguyen Tien Dat\textsuperscript{1,\,*}, Nguyen Diep Linh\textsuperscript{1}, Bach Duc Anh\textsuperscript{1}, Bui Thi Thu Thuy\textsuperscript{1} and Pham Le Nhat Linh\textsuperscript{1}}
\end{center}

\begin{center}
{\fontsize{11}{13}\itshape\selectfont
\textsuperscript{1}Faculty of Information Technology, School of Mathematics and Information Technology,\\
Hanoi National University of Education, Hanoi, Vietnam}
\end{center}

\begin{center}
{\fontsize{11}{13}\selectfont
\textsuperscript{*}Corresponding author: Nguyen Tien Dat, e-mail: \href{mailto:mas835151465@hnue.edu.vn}{mas835151465@hnue.edu.vn} (primary); \href{mailto:datlcpro@gmail.com}{datlcpro@gmail.com} (alternative)}\\[1pt]
{\fontsize{10}{12}\selectfont
Received \today. Revised TBD. Accepted TBD.}
\end{center}

\vspace{4pt}

{\fontsize{10}{12}\selectfont
\setlength{\leftskip}{1cm}\setlength{\rightskip}{1cm}
\noindent\textbf{Abstract.} Item Response Theory (IRT) systems deployed at scale, including Vietnam's OLM platform, typically assume a static student ability~\(\theta\), yet large-scale online learning produces densely longitudinal data that contradicts this assumption. We construct and validate an integrated IRT--LSEM framework for estimating dynamic ability~\(\theta_{it}\) on real OLM Mathematics data covering 24{,}213 high-school students and 122{,}957 assessment occasions across grades 10--12. The framework couples a 1PL measurement layer (EAP scoring) with four structural layers (LGCM, LCSM, multilevel growth, continuous-time DSEM) and a Kalman smoother that exploits EAP standard errors as time-varying observation noise. The dynamic model beats the static-Rasch baseline by \(\Delta\)AIC = 82--167 (LRT \(p \le 10^{-19}\)); LCSM reveals catch-up at grade 10 (\(\beta = -0.069\), CI excludes 0) shifting to plateau at grade 12; Kalman cuts mean EAP error variance by 32--41\% with no extra items.\par
\noindent\textbf{\textit{Keywords:}} \textit{Item Response Theory, longitudinal structural equation modeling, Kalman smoothing, dynamic ability estimation, educational data mining.}
\par}

\vspace{6pt}
